\section{Introduction}\label{sec:introduction}

\begin{figure}[b!]
  \centering
  \includegraphics[width=\linewidth]{fig1-overview.pdf}
  \caption{Results overview. (a) Average scores on four streaming
  benchmarks against open-source streaming baselines (details in
  \Cref{tab:streaming}; the StreamingBench average covers its visual
  groups, \S\ref{sec:eval-streaming}). (b) \mossvl{} against the best competing model on
  the proactive streaming subsets---PA (Proactive Alerting, OmniMMI),
  PO (Proactive Output, StreamingBench), and FAR (Forward Active
  Responding, OVO-Bench), all testing whether the model speaks up
  unprompted at the right moment---and on selected offline strengths
  (details in \Cref{tab:streaming,tab:offline}). Green = \mossvl{}; gray =
  the labeled competitor.}
  \label{fig:overview}
\end{figure}

\begin{table}[t]
  \centering
  \caption{Levels of video understanding, from offline to real-time. Each
  level lights up one additional capability axis; the dividing line between
  the streaming regime (L2--L4) and real-time (L5) is whether the model
  keeps perceiving \emph{while} it generates---L2--L4 models are blind
  during a reply, an L5 model is not. \mossrealtime{} achieves L5 behavior,
  demonstrated through live demos and the released real-time inference
  code, and is quantitatively validated at levels L2--L4 on four streaming
  benchmarks; a dedicated benchmark for L5 behavior remains an open
  problem.}
  \label{tab:levels}
  \small
  \renewcommand{\arraystretch}{1.3}
  \begin{tabular}{@{}clp{5.4cm}cccc@{}}
    \toprule
    & & & \multicolumn{4}{c}{Capability axes} \\
    \cmidrule(lr){4-7}
    Level & Regime & Defining behavior
      & \makecell{Stream.\\input}
      & \makecell{Silence}
      & \makecell{Multi-\\resp.}
      & \makecell{Perceive\\while gen.} \\
    \midrule
    L1 & offline & Watches the full video first, then answers questions
         about it (multi-turn allowed). & \xmark & -- & \xmark & \xmark \\
    L2 & streaming & Video arrives continuously; the user may ask at any
         moment and the model answers immediately---but it is blind while
         replying. & \cmark & \xmark & \xmark & \xmark \\
    L3 & streaming & Adds \emph{waiting}: if the answer is not yet
         determinable, the model stays silent until the key evidence
         appears, then answers. & \cmark & \cmark & \xmark & \xmark \\
    L4 & streaming & Adds \emph{persistent queries}: one question stays
         resident and the model re-answers as the scene evolves (still
         blind during each individual reply). & \cmark & \cmark & \cmark
         & \xmark \\
    L5 & real-time & Adds \emph{perception during generation}: the model
         keeps watching while it speaks, revising or interrupting its own
         reply the moment the evidence changes. & \cmark & \cmark & \cmark
         & \cmark \\
    \bottomrule
  \end{tabular}
\end{table}


Most open vision--language models understand video offline: given a
finished clip, they read it end to end and then answer questions about
it \citep{Bai2025qwen,An2026llava,Team2026gemma}. The settings where
video understanding matters most do not wait for the clip to end. A
live assistant watches a scene that is still unfolding, decides for
itself when something is worth saying, and must keep watching while it
says it. Table~\ref{tab:levels} organizes this capability space into
five levels. L1 is the offline regime. L2--L4 form the streaming
regime occupied by recent streaming models
\citep{Wang2024videollm,Fu2025vispeak,Xia2025streaming,Lu2026aura}:
input arrives continuously, deliberate silence and persistent queries
come into play, yet the model stays blind for the duration of each
reply. L5 adds the ability that separates real-time interaction from
everything below: perceiving \emph{while} generating, so a reply can be
revised or cut short the moment the evidence changes.

\mossvl{} is an open vision--language model family that treats
real-time interaction as a first-class capability, and it reaches that
capability by co-design rather than through any single component. The
architecture enables the behavior: the language decoder attends to
vision only through gated cross-attention, so visual tokens never enter
the decoded sequence, and an arriving frame merely appends to the
cross-attention cache---the model naturally keeps perceiving while it
generates (\S\ref{sec:architecture}). XRoPE orders text and vision along one
shared timeline, and absolute timestamp tokens make wall-clock time
explicit. The data injects the behavior: a synthesized corpus of
real-time interaction supervises when to speak, when to stay silent,
and how to revise a reply the scene has overturned
(\S\ref{sec:posttraining}). The training strategy keeps the stack
stable: a four-stage pre-training curriculum and standard supervised
fine-tuning build the offline foundation (\S\ref{sec:pretraining}),
every real-time-specific choice is concentrated in Realtime-SFT, one
light final stage, and a single system prompt moves the same weights
among offline, streaming, and real-time operation. \mossvl{} continues
a line that began with MOSS-Video-Preview \citep{Wang2026moss}, which
explored the real-time paradigm; this release redesigns the stack from
scratch and gives the line its first quantitative streaming evaluation.

Figure~\ref{fig:overview} previews the outcome. Offline,
\mossinstruct{} is competitive with open models of comparable scale
and leads the temporal-reasoning video sets Minerva, TOMATO, and
VideoMME-Logical (\S\ref{sec:eval-offline}). In the streaming regime,
the wins land precisely where timing is being tested: across four
streaming benchmarks against open-source streaming baselines,
\mossrealtime{} posts the best average on three of the four (and is
second on the fourth), sweeping the three subsets that squarely test
proactive behavior---66.0 vs.\ 37.5 on OmniMMI's Proactive Alerting
(\S\ref{sec:eval-streaming}). Efficiency follows from the same design:
against Qwen3-VL-8B, built on the same Qwen3-8B language backbone, the
time-to-first-token gap widens from $2.8\times$ to $5.1\times$ as
visual context grows (\S\ref{sec:eval-efficiency}). L5 behavior itself
is demonstrated qualitatively, through live demos and the released
real-time inference code (\S\ref{sec:eval-showcase}); quantitative
validation covers L2--L4, where public benchmarks exist.

\pagebreak
In summary, our main contributions are:
\begin{itemize}[leftmargin=1.5em, itemsep=2pt, topsep=2pt]
  \item \textbf{A stack co-designed for real-time interaction.} The
  architecture enables it---gated cross-attention with XRoPE and
  absolute timestamps lets the model take in new frames naturally
  while it generates;
  the data injects it---synthesized streams supervise response
  timing; the training strategy keeps it stable---the language backbone
  stays intact behind zero-initialized gates.
  \item \textbf{Realtime-SFT, an interaction paradigm in one light
  stage.} Two added state tokens, one shared system prompt, and a
  reweighted next-token loss---under 3\% of total training
  tokens---teach when to speak, when to stay silent, and when to
  revise, delivering the best average on three of four streaming
  benchmarks and the proactive sweep above.
  \item \textbf{Timing understanding shows up offline as well.}
  \mossinstruct{}, trained without the real-time corpus, leads the
  temporal-reasoning video sets in our
  comparison---Minerva, TOMATO, VideoMME-Logical---echoing the
  perception and temporal-understanding foundations laid during
  pre-training.
  \item \textbf{More parameters, faster serving.} \mossvl{} spends 11.3B
  parameters, yet visual tokens stay outside the decoded sequence, so
  serving latency grows more slowly with visual context than its
  same-backbone interleaved counterpart---measured in SGLang, not
  estimated. We release the complete training curriculum and all five
  checkpoints.
\end{itemize}
